\section{Sharpen the continuous C wedge using reference numerator mass}
The prior unbounded factor-count wedge is an input. This refinement
strictly enlarges the ALL-padding actual late-C wedge to
\[
 4(u+r)+17>=13(h-u),\quad\hbox{equivalently }17u+4r+17>=13h.
\]
Only delta=h-u>=5 is new. Put w=u+r,W=w+delta,p=delta-2.
The new cone implies W>=17 and w+1>=(13/4)(delta-1).
The registered CI-count loss theorem therefore gives
l_j>=(9/4)l_(j-1) for all1<=j<=p+1. In particular earlier losses
double, sum_(j<p)l_j<=b=l_p, and c=l_(p+1)>=9b/4.

\paragraph{An all-word reference numerator-mass bound.}
Let B be the nonnegative reference numerator, M=B(1), P=prod(V).
Its exact proper suffix expansion gives M=P-A_V(1), hence0<M<P.
For t>=3 let F_t=phi^t S_t(q)=phi^(t+1)-phi.
For EVERY t>=3,
\[
 F_t>=t(8/7)^t.
\]
The base t3 uses F3=2phi+2>=26/5>3(8/7)^3. The recurrence
F_(t+1)=phi F_t+1 and phi*t>=(8/7)(t+1) for ALLt>=3
prove the induction. Multiply over ALL factors:
B(q)phi^W>=P(8/7)^W. In the Binet formula the conjugate term
in x=Dref_W is POSITIVE because sign B(-phi)=(-1)^W. Therefore
\[
 x>=P(8/7)^W/(phi sqrt5)>=2P\quad(W>=15).
\]
The last inequality uses (8/7)^15>=117/16,
phi<=13/8,sqrt5<=9/4. Thus the ENTIRE numerator mass obeys M<=x/2.

\paragraph{A stronger homogeneous reference ratio and exact cumulative reads.}
Whole-word spectral norm gives t<=q^(W+r-2)<=q^10=34-55q
at W+r>=12. Then y/x=(phi-qt)/(1+t)>=8/5, with equality at q^10.
The coefficient identity
\[
 g_n=R_(n+2)-sum_(i<=n+2)B_i
\]
follows from g's binary kernel Fib(n+3)-1. Consequently
\[
 g_(W-2)=x+y-M>=2x,\quad
 g_(W-1)=x+2y-M>=7x/2.
\]
All these bounds apply in the new cone since W>=17.
No global reference cumulative LC is used.

\paragraph{Direct payment of every cumulative read with only doubling losses.}
As in the first wedge, consider the strictly growing actual pair
D_W=x-b,D_(W+1)=y-c and lambda=(y-c)/(x-b)>1.
Because c>=9b/4 and phi<=13/8, the exact entry identity gives
\[
 lambda<phi,\qquad(phi-lambda)x>5b/8.
\]
For a read k=W-d inside the loss band, the actual cumulative source
comparison is at most S_(p-d)<=2l_(p-d)<=2b/2^d.
If d>=1, reference dyadic upper growth and the stronger mass bound give
g_(k-1)>=g_(W-2)/2^(d-1)>=4x/2^d. Its residual is therefore
strictly greater than (5/2)b/2^d, paying the entire source.
At d0 use g_(W-1)>=7x/2, giving residual>35b/16>2b.
Earlier untouched cumulative ratios still exceed phi. Every boundary
G comparison through W is paid, with a strict margin.

\paragraph{The ordinary seed works with a<=b, not just a<=b/2.}
Put a=sum_(j<p)l_j, so0<=a<=b. The exact ordinary determinant is
Q=A-x(a+2b)+y(a-b+c)+b^2+ab-ac, A=x^2+xy-y^2>0.
For ANY reals x>=b>0,y>=8x/5,0<=a<=b,c>=9b/4, one has
\[
 Q>=A+7b^2/20>0.
\]
To prove it, first replace y by8x/5 (its coefficient is positive).
The coefficient of x is (3/5)a-(18/5)b+(8/5)c>=0,
so replace x by b. The coefficient of c becomes (8/5)b-a>=0,
so replace c by9b/4. The result is A+b^2-(13/20)ab,
which is at least A+7b^2/20. This is a generic unbounded cone proof.

Universal late-D-LC plus the paid G prefix and ordinary seed now
integrates every later all-shift C comparison. No new finite entrance
is needed. Earlier C, A2 and original UI remain OPEN. The unbounded interval-factor amplification and finite-sum loss
inductions are kernel checked, along with complete entry and read
payments. The reference/model/product/source identities and read
integration are written assemblies. No end-to-end Lean composition
theorem is claimed, and no finite entrance premise is introduced.

\paragraph{Statement and scope.}
For EVERY actual sorted word $1^\alpha2^uV$, with nonempty capacities
$V\ge3$, $h=\sum(t-1)$ and $r=|V|$, the condition
$17u+4r+17\ge13h$ implies EVERY C shift at least three at EVERY
current rank $N\ge W+1$, for EVERY padding. Degree gaps at most four
use the prior registered theorem; the ENTIRE gap-at-least-five cone
is unbounded and analytic here. Potential late-C failures now also
require $17u+4r+17<13h$, in addition to the existing residual bounds.
General earlier C, A2 and original coefficient unimodality remain OPEN.
